\documentclass[11pt]{article}

\usepackage[margin=1in]{geometry}
\usepackage{amsmath,amssymb,amsthm}

\theoremstyle{definition}
\newtheorem*{solution}{Solution}

\newcommand{\A}{\mathbb{A}}
\newcommand{\C}{\mathbb{C}}
\newcommand{\Pj}{\mathbb{P}}
\newcommand{\Q}{\mathbb{Q}}
\newcommand{\Z}{\mathbb{Z}}
\newcommand{\K}{K}
\newcommand{\V}{\operatorname{V}}
\newcommand{\Frac}{\operatorname{Frac}}
\newcommand{\Gal}{\operatorname{Gal}}

\title{Worked Solutions Companion\\
       \large Algebraic Geometry and Galois Theory}
\author{Companion to \texttt{algebraic\_geometry\_notes\_2026\_06\_06.tex}}
\date{June 6, 2026}

\begin{document}

\maketitle
\tableofcontents

\section*{About This Companion}
\addcontentsline{toc}{section}{About This Companion}

This document gives worked solutions to the exercises in the lecture chapter
\emph{Algebraic Geometry and Galois Theory}.  The solutions are standard
background calculations written to be read alongside the notes.  They are not
additional recovered content from the handwritten source PDF.

\section{Field Extensions}

\subsection*{Exercise 12.1}
\addcontentsline{toc}{subsection}{Exercise 12.1}
\textbf{Problem.}  Let \(F\subseteq E\) be a finite field extension.  Explain
why \(E\) is an \(F\)-vector space, and identify scalar multiplication.

\begin{solution}
The elements of the vector space are the elements of the larger field \(E\).
The addition operation is addition in the field \(E\).  Since \(F\subseteq E\),
each scalar \(a\in F\) is also an element of \(E\), so scalar multiplication is
field multiplication:
\[
  F\times E\longrightarrow E,
  \qquad
  (a,v)\longmapsto av.
\]
The vector-space axioms follow from the field axioms in \(E\).  For example,
if \(a,b\in F\) and \(u,v\in E\), then
\[
  a(u+v)=au+av,
  \qquad
  (a+b)v=av+bv,
  \qquad
  (ab)v=a(bv),
  \qquad
  1_Fv=v.
\]
The extension is finite exactly when this vector space has finite dimension:
\[
  [E:F]=\dim_F E<\infty.
\]
\end{solution}

\subsection*{Exercise 12.2}
\addcontentsline{toc}{subsection}{Exercise 12.2}
\textbf{Problem.}  For \(\Q(\sqrt d)/\Q\), with \(d\) squarefree and not a
square, write the matrix for multiplication by \(a+b\sqrt d\) in the basis
\((1,\sqrt d)\).

\begin{solution}
Every element of \(\Q(\sqrt d)\) has a unique form
\[
  r+s\sqrt d,\qquad r,s\in\Q.
\]
Let \(u=a+b\sqrt d\).  Multiplication by \(u\) is the \(\Q\)-linear map
\[
  m_u\colon \Q(\sqrt d)\longrightarrow \Q(\sqrt d),
  \qquad
  v\longmapsto uv.
\]
Compute
\[
  (a+b\sqrt d)(r+s\sqrt d)
  =
  ar+as\sqrt d+br\sqrt d+bs d
  =
  (ar+bds)+(as+br)\sqrt d.
\]
Thus the coordinate column of \(m_u(r+s\sqrt d)\) in the basis
\((1,\sqrt d)\) is
\[
  \begin{pmatrix}
    ar+bds\\
    br+as
  \end{pmatrix}
  =
  \begin{pmatrix}
    a & bd\\
    b & a
  \end{pmatrix}
  \begin{pmatrix}
    r\\
    s
  \end{pmatrix}.
\]
Therefore
\[
  [m_{a+b\sqrt d}]_{(1,\sqrt d)}
  =
  \begin{pmatrix}
    a & bd\\
    b & a
  \end{pmatrix}.
\]
\end{solution}

\section{Coordinate Rings and Function Fields}

\subsection*{Exercise 12.3}
\addcontentsline{toc}{subsection}{Exercise 12.3}
\textbf{Problem.}  Compute the coordinate ring and function field of
\[
  X=\V(y-x^2)\subseteq \A^2_{\C}.
\]

\begin{solution}
By definition,
\[
  \C[X]=\C[x,y]/(y-x^2).
\]
Define
\[
  \Phi\colon \C[x,y]\longrightarrow \C[t],
  \qquad
  \Phi(x)=t,\quad \Phi(y)=t^2.
\]
Then
\[
  \Phi(y-x^2)=t^2-t^2=0,
\]
so \((y-x^2)\subseteq\ker\Phi\).

Conversely, reduce any polynomial \(f(x,y)\) modulo the relation
\(y=x^2\).  Since \(y-x^2\) is monic in \(y\), division by \(y-x^2\) gives
\[
  f(x,y)=q(x,y)(y-x^2)+r(x),
\]
where \(r(x)=f(x,x^2)\).  If \(f\in\ker\Phi\), then
\[
  0=\Phi(f)=f(t,t^2)=r(t).
\]
Hence \(r=0\), so \(f\in (y-x^2)\).  Therefore
\[
  \ker\Phi=(y-x^2).
\]
By the first isomorphism theorem,
\[
  \C[x,y]/(y-x^2)\cong \C[t].
\]
The function field is the fraction field:
\[
  \K(X)=\Frac(\C[X])\cong \Frac(\C[t])=\C(t).
\]
\end{solution}

\subsection*{Exercise 12.4}
\addcontentsline{toc}{subsection}{Exercise 12.4}
\textbf{Problem.}  Compute the coordinate ring and function field of
\[
  U=\V(xy-1)\subseteq \A^2_{\C}.
\]

\begin{solution}
The coordinate ring is
\[
  \C[U]=\C[x,y]/(xy-1).
\]
In this quotient, the relation \(xy-1=0\) says
\[
  xy=1.
\]
Thus the class of \(y\) is the inverse of the class of \(x\):
\[
  y=x^{-1}.
\]

To express this as an isomorphism, define
\[
  \Psi\colon \C[x,y]\longrightarrow \C[t,t^{-1}],
  \qquad
  \Psi(x)=t,\quad \Psi(y)=t^{-1}.
\]
Then
\[
  \Psi(xy-1)=tt^{-1}-1=0,
\]
so \((xy-1)\subseteq\ker\Psi\).  The map is surjective because \(t=\Psi(x)\)
and \(t^{-1}=\Psi(y)\).  Reducing by the relation \(xy=1\) lets every element
be written using powers of \(x\) and \(x^{-1}\), so the kernel is exactly
\((xy-1)\).  Hence
\[
  \C[x,y]/(xy-1)\cong \C[t,t^{-1}].
\]
Therefore
\[
  \K(U)=\Frac(\C[t,t^{-1}])=\C(t).
\]
\end{solution}

\section{Projective Coordinates}

\subsection*{Exercise 12.5}
\addcontentsline{toc}{subsection}{Exercise 12.5}
\textbf{Problem.}  Show that \(XZ-Y^2=0\) is well-defined on
\(\Pj^2_{\C}\).

\begin{solution}
In projective space,
\[
  [X:Y:Z]=[\lambda X:\lambda Y:\lambda Z]
  \qquad
  \text{for }\lambda\in\C^\times.
\]
Let
\[
  F(X,Y,Z)=XZ-Y^2.
\]
Then
\[
  F(\lambda X,\lambda Y,\lambda Z)
  =
  (\lambda X)(\lambda Z)-(\lambda Y)^2
  =
  \lambda^2(XZ-Y^2)
  =
  \lambda^2F(X,Y,Z).
\]
Since \(\lambda^2\neq 0\), we have
\[
  F(\lambda X,\lambda Y,\lambda Z)=0
  \quad\Longleftrightarrow\quad
  F(X,Y,Z)=0.
\]
Thus the truth of \(XZ-Y^2=0\) does not depend on the representative of the
projective point.
\end{solution}

\subsection*{Exercise 12.6}
\addcontentsline{toc}{subsection}{Exercise 12.6}
\textbf{Problem.}  Give an example showing that \(X+1=0\) is not
well-defined on \(\Pj^1_{\C}\).

\begin{solution}
The two nonzero vectors \((-1,1)\) and \((-2,2)\) define the same point of
\(\Pj^1_{\C}\), because
\[
  (-2,2)=2(-1,1).
\]
Thus
\[
  [-1:1]=[-2:2].
\]
Now test the equation \(X+1=0\) on the two representatives:
\[
  -1+1=0,
  \qquad
  -2+1=-1\neq 0.
\]
The same projective point satisfies the equation for one representative and
does not satisfy it for another.  Therefore \(X+1=0\) is not a well-defined
equation on projective space.
\end{solution}

\section{Veronese Maps}

\subsection*{Exercise 12.7}
\addcontentsline{toc}{subsection}{Exercise 12.7}
\textbf{Problem.}  For
\[
  \nu_2([x:y])=[x^2:xy:y^2],
\]
verify \(XZ-Y^2=0\) and recover the parameter on \(X\neq 0\).

\begin{solution}
Let
\[
  [X:Y:Z]=[x^2:xy:y^2]
\]
be a point in the image.  Then
\[
  XZ-Y^2=(x^2)(y^2)-(xy)^2=x^2y^2-x^2y^2=0.
\]
Thus every image point satisfies the conic equation.

On the target chart \(X\neq 0\), the source coordinate \(x\) is nonzero.  We
may write the source point as
\[
  [x:y]=[1:t],
  \qquad
  t=y/x.
\]
Then
\[
  \nu_2([1:t])=[1:t:t^2].
\]
In target coordinates, this means
\[
  t=\frac{Y}{X},
  \qquad
  \frac{Z}{X}=t^2.
\]
So the parameter is recovered by \(t=Y/X\) on the chart \(X\neq 0\).
\end{solution}

\subsection*{Exercise 12.8}
\addcontentsline{toc}{subsection}{Exercise 12.8}
\textbf{Problem.}  Repeat the injectivity argument for
\[
  \nu_3([x:y])=[x^3:x^2y:xy^2:y^3].
\]

\begin{solution}
Use source coordinates \([u:v]\) to avoid confusing them with target
coordinates.  The map is
\[
  \nu_3([u:v])=[u^3:u^2v:uv^2:v^3].
\]
If \(u\neq 0\), then
\[
  [u:v]=[1:t],
  \qquad
  t=v/u.
\]
The image is
\[
  \nu_3([1:t])=[1:t:t^2:t^3].
\]
Thus on the target chart \(X\neq 0\), the parameter is recovered as
\[
  t=\frac{Y}{X}.
\]
Hence the image determines the source point \([1:t]\).

If \(u=0\), then the source point is \([0:1]\), and
\[
  \nu_3([0:1])=[0:0:0:1].
\]
This target point cannot arise from any source point with \(u\neq 0\), because
those have \(X\neq 0\).  Therefore each image point has a unique preimage, so
\(\nu_3\) is injective.
\end{solution}

\subsection*{Exercise 12.9}
\addcontentsline{toc}{subsection}{Exercise 12.9}
\textbf{Problem.}  Verify the three equations
\[
  Y^2=XZ,\qquad Z^2=YW,\qquad YZ=XW
\]
on the image of \(\nu_3\).

\begin{solution}
For a point in the image of \(\nu_3\), write
\[
  X=x^3,\qquad
  Y=x^2y,\qquad
  Z=xy^2,\qquad
  W=y^3.
\]
Then
\[
  Y^2=(x^2y)^2=x^4y^2
  \qquad\text{and}\qquad
  XZ=(x^3)(xy^2)=x^4y^2,
\]
so \(Y^2=XZ\).  Also
\[
  Z^2=(xy^2)^2=x^2y^4
  \qquad\text{and}\qquad
  YW=(x^2y)(y^3)=x^2y^4,
\]
so \(Z^2=YW\).  Finally,
\[
  YZ=(x^2y)(xy^2)=x^3y^3
  \qquad\text{and}\qquad
  XW=(x^3)(y^3)=x^3y^3,
\]
so \(YZ=XW\).  Hence every image point satisfies all three equations.
\end{solution}

\section{The Elliptic Curve}

\subsection*{Exercise 12.10}
\addcontentsline{toc}{subsection}{Exercise 12.10}
\textbf{Problem.}  On the chart \(Z\neq 0\), reduce
\[
  Y^2Z=X^3+aXZ^2+bZ^3
\]
to \(y^2=x^3+ax+b\).  Then compute the point with \(Z=0\).

\begin{solution}
On \(Z\neq 0\), define affine coordinates
\[
  x=X/Z,\qquad y=Y/Z.
\]
Divide the homogeneous equation by \(Z^3\):
\[
  \frac{Y^2Z}{Z^3}
  =
  \frac{X^3}{Z^3}
  +a\frac{XZ^2}{Z^3}
  +b\frac{Z^3}{Z^3}.
\]
This gives
\[
  \left(\frac{Y}{Z}\right)^2
  =
  \left(\frac{X}{Z}\right)^3
  +a\left(\frac{X}{Z}\right)+b,
\]
or
\[
  y^2=x^3+ax+b.
\]

Now set \(Z=0\) in the homogeneous equation.  We get
\[
  0=X^3,
\]
so \(X=0\).  Since \([0:0:0]\) is not a projective point, we must have
\(Y\neq 0\).  Rescale by \(Y^{-1}\):
\[
  [0:Y:0]=[0:1:0].
\]
Thus the only point of the curve with \(Z=0\) is \([0:1:0]\).
\end{solution}

\subsection*{Exercise 12.11}
\addcontentsline{toc}{subsection}{Exercise 12.11}
\textbf{Problem.}  Verify that \((x,y)\mapsto (x,-y)\) preserves
\[
  y^2=x^3+ax+b.
\]

\begin{solution}
Suppose \((x,y)\) lies on the affine curve.  Then
\[
  y^2=x^3+ax+b.
\]
For the transformed point \((x,-y)\), the left side of the equation becomes
\[
  (-y)^2=y^2.
\]
Therefore
\[
  (-y)^2=y^2=x^3+ax+b.
\]
So \((x,-y)\) also lies on the same affine curve.  Hence the transformation
preserves the equation.
\end{solution}

\section{Function Fields, Invariants, and Galois Theory}

\subsection*{Exercise 12.12}
\addcontentsline{toc}{subsection}{Exercise 12.12}
\textbf{Problem.}  In \(\C(x)(y)\), where \(y^2=f(x)=x^3+ax+b\), reduce
\[
  c_0(x)+c_1(x)y+c_2(x)y^2+c_3(x)y^3
\]
to \(A(x)+B(x)y\).

\begin{solution}
Use the defining relation
\[
  y^2=f(x).
\]
Then
\[
  y^3=y\cdot y^2=yf(x)=f(x)y.
\]
Therefore
\[
\begin{aligned}
  &c_0(x)+c_1(x)y+c_2(x)y^2+c_3(x)y^3 \\
  &\qquad
  =
  c_0(x)+c_1(x)y+c_2(x)f(x)+c_3(x)f(x)y\\
  &\qquad
  =
  \bigl(c_0(x)+c_2(x)f(x)\bigr)
  +
  \bigl(c_1(x)+c_3(x)f(x)\bigr)y.
\end{aligned}
\]
Thus
\[
  A(x)=c_0(x)+c_2(x)f(x),
  \qquad
  B(x)=c_1(x)+c_3(x)f(x).
\]
\end{solution}

\subsection*{Exercise 12.13}
\addcontentsline{toc}{subsection}{Exercise 12.13}
\textbf{Problem.}  Using
\[
  \K(E)=\C(x)\oplus \C(x)y,
\]
show that
\[
  \K(E)^{\langle\iota^*\rangle}=\C(x),
\]
where \(\iota^*(x)=x\) and \(\iota^*(y)=-y\).

\begin{solution}
Let \(u\in\K(E)\).  Using the vector-space decomposition, write
\[
  u=a(x)+b(x)y,
  \qquad
  a(x),b(x)\in \C(x).
\]
Since \(\iota^*(x)=x\), every rational function of \(x\) is fixed:
\[
  \iota^*(a(x))=a(x),
  \qquad
  \iota^*(b(x))=b(x).
\]
Since \(\iota^*(y)=-y\), we get
\[
  \iota^*(u)
  =
  a(x)-b(x)y.
\]
The element \(u\) is invariant if and only if
\[
  a(x)+b(x)y=a(x)-b(x)y.
\]
Subtracting gives
\[
  2b(x)y=0.
\]
The field has characteristic zero and \(y\neq 0\), so \(b(x)=0\).  Hence every
invariant element is \(a(x)\in \C(x)\).  Conversely, every element of
\(\C(x)\) is fixed because \(\iota^*(x)=x\).  Therefore
\[
  \K(E)^{\langle\iota^*\rangle}=\C(x).
\]
\end{solution}

\subsection*{Exercise 12.14}
\addcontentsline{toc}{subsection}{Exercise 12.14}
\textbf{Problem.}  Let \(L=\C(x)(y)\), where \(y^2=x^3+ax+b\).  Verify that
\[
  \sigma(x)=x,\qquad \sigma(y)=-y
\]
defines a nontrivial \(\C(x)\)-automorphism of \(L\), and compute
\(\sigma^2\).

\begin{solution}
Let
\[
  f(x)=x^3+ax+b.
\]
The field \(L\) is generated over \(\C(x)\) by \(y\), with relation
\[
  y^2=f(x).
\]
The proposed map fixes every element of \(\C(x)\), because it fixes \(x\).
It sends \(y\) to \(-y\).  This preserves the defining relation:
\[
  \sigma(y)^2=(-y)^2=y^2=f(x)=\sigma(f(x)).
\]
Therefore the assignment extends to a \(\C(x)\)-automorphism of \(L\).

It is nontrivial because
\[
  \sigma(y)=-y\neq y
\]
in the function field.  Applying \(\sigma\) twice gives
\[
  \sigma^2(x)=\sigma(x)=x,
  \qquad
  \sigma^2(y)=\sigma(-y)=-\sigma(y)=-(-y)=y.
\]
Since \(L=\C(x)(y)\) is generated by \(x\) and \(y\), this proves
\[
  \sigma^2=\operatorname{id}_L.
\]
Thus the automorphism has order two.
\end{solution}

\subsection*{Exercise 12.15}
\addcontentsline{toc}{subsection}{Exercise 12.15}
\textbf{Problem.}  Explain why invariant rational functions are constant on
paired points \(P\) and \(\iota(P)\), and why they descend to the quotient.

\begin{solution}
Let \(f\in\K(E)\) be invariant under the involution.  Algebraically this means
\[
  \iota^*(f)=f.
\]
The pullback \(\iota^*(f)\) is the rational function obtained by first applying
\(\iota\) and then evaluating \(f\).  Thus, wherever both sides are defined,
\[
  \iota^*(f)(P)=f(\iota(P)).
\]
The equality \(\iota^*(f)=f\) therefore gives
\[
  f(P)=f(\iota(P)).
\]
So \(f\) has the same value on the paired points \(P\) and \(\iota(P)\).

The quotient \(E/\langle\iota\rangle\) identifies \(P\) and \(\iota(P)\).
Because \(f\) takes the same value on the two representatives of each generic
equivalence class, its value depends only on the quotient point.  Hence \(f\)
descends to a rational function on the quotient.  In the elliptic-curve
example, the quotient has function field \(\C(x)\), so the invariant rational
functions are exactly rational functions of the \(x\)-coordinate.
\end{solution}

\end{document}
